\begin{derivation}[从\eqnrefs{eq:sm-heavy-block-schur-matching-dp85}到\eqnrefs{eq:sm-heavy-field-local-control-dp85}]
把轻、重自由度组成列向量
\begin{equation*}
 \Psi\equiv
 \begin{pmatrix}
  \psi_L\\ \Psi_H
 \end{pmatrix},
\end{equation*}
并把二次型作用量写成
\begin{equation*}
 S_2
 =\frac12\int_p
 \Psi^\dagger(-p)
 \begin{pmatrix}
 K_{LL}(p)&K_{LH}(p)\\
 K_{HL}(p)&K_{HH}(p)
 \end{pmatrix}
 \Psi(p).
\end{equation*}
相应的线性运动方程是
\begin{align*}
 K_{LL}\psi_L+K_{LH}\Psi_H&=0,\\
 K_{HL}\psi_L+K_{HH}\Psi_H&=0.
\end{align*}
只要过程能区内 $K_{HH}(p)$ 可逆，第二式可以精确解为
\begin{equation*}
 \Psi_H=-K_{HH}^{-1}K_{HL}\psi_L.
\end{equation*}
代回第一式：
\begin{align*}
 0
 &=K_{LL}\psi_L+K_{LH}\Psi_H\\
 &=K_{LL}\psi_L
 -K_{LH}K_{HH}^{-1}K_{HL}\psi_L\\
 &=\left(K_{LL}-K_{LH}K_{HH}^{-1}K_{HL}\right)\psi_L.
\end{align*}
因此轻场的精确二次核为
\begin{equation*}
 K_{\rm eff}(p)
 =K_{LL}(p)-K_{LH}(p)K_{HH}^{-1}(p)K_{HL}(p),
\end{equation*}
即正文\eqnrefs{eq:sm-heavy-block-schur-matching-dp85}。

同一个结果也可由分块高斯消元直接验证。定义下三角矩阵
\begin{equation*}
 L=
 \begin{pmatrix}
 I&-K_{LH}K_{HH}^{-1}\\
 0&I
 \end{pmatrix}.
\end{equation*}
则
\begin{align*}
 L
 \begin{pmatrix}
 K_{LL}&K_{LH}\\
 K_{HL}&K_{HH}
 \end{pmatrix}
 &=
 \begin{pmatrix}
 K_{LL}-K_{LH}K_{HH}^{-1}K_{HL}&0\\
 K_{HL}&K_{HH}
 \end{pmatrix}.
\end{align*}
所以 Schur 补不是低能近似，而是对重块可逆这一条件下的精确代数恒等式。

要得到局域 EFT，才需要进一步使用重--轻尺度分离。把重块拆成
\begin{equation*}
 K_{HH}(p)=K_M+\delta K(p),
\end{equation*}
其中 $K_M$ 是由重质量主导、在 $p=0$ 已可逆的部分。提出 $K_M$：
\begin{align*}
 K_{HH}(p)
 &=K_M\left[I+K_M^{-1}\delta K(p)\right].
\end{align*}
定义无量纲算符
\begin{equation*}
 X(p)\equiv K_M^{-1}\delta K(p).
\end{equation*}
若在目标能区内
\begin{equation*}
 \epsilon_{\rm Schur}\equiv\sup_{p\in\mathcal D}\|X(p)\|<1,
\end{equation*}
则 Neumann 级数在算符范数下绝对收敛：
\begin{align*}
 K_{HH}^{-1}(p)
 &=[I+X(p)]^{-1}K_M^{-1}\\
 &=\sum_{n=0}^{\infty}[-X(p)]^nK_M^{-1}.
\end{align*}
截断到 $N$ 阶时，利用有限几何级数恒等式
\begin{equation*}
 [I+X]^{-1}
 =\sum_{n=0}^{N}(-X)^n
 +(-X)^{N+1}[I+X]^{-1},
\end{equation*}
得到预解式余项
\begin{equation*}
 R_{N+1}^{(H)}(p)
 =(-X)^{N+1}[I+X]^{-1}K_M^{-1}.
\end{equation*}
由次乘性和
\begin{equation*}
 \|[I+X]^{-1}\|
 \le\sum_{r=0}^{\infty}\|X\|^r
 \le\frac1{1-\epsilon_{\rm Schur}}
\end{equation*}
可得严格范数界
\begin{equation*}
 \|R_{N+1}^{(H)}(p)\|
 \le
 \|K_M^{-1}\|
 \frac{\epsilon_{\rm Schur}^{N+1}}{1-\epsilon_{\rm Schur}}.
\end{equation*}
代回 Schur 补以后，把预解式级数严格截断到 $X^N$，局域 EFT 二次核为
\begin{align*}
 K_{\rm eff}^{(N)}
 &=K_{LL}
 -K_{LH}
 \left[\sum_{n=0}^{N}(-X)^nK_M^{-1}\right]
 K_{HL}\\
 &=K_{LL}
 -K_{LH}K_M^{-1}K_{HL}
 +K_{LH}XK_M^{-1}K_{HL}\\
 &\quad-K_{LH}X^2K_M^{-1}K_{HL}
 +\cdots
 +(-1)^{N+1}K_{LH}X^NK_M^{-1}K_{HL}.
\end{align*}
这里 $X=K_M^{-1}\delta K$，所以每增加一个 $X$ 就增加一个轻尺度相对于重尺度的幂。精确核与该有限和之差完全由上面的 $R_{N+1}^{(H)}$ 给出，并满足
\begin{equation*}
 \|K_{\rm eff}-K_{\rm eff}^{(N)}\|
 \le
 \|K_{LH}\|\,\|K_{HL}\|\,
 \|K_M^{-1}\|
 \frac{\epsilon_{\rm Schur}^{N+1}}{1-\epsilon_{\rm Schur}}.
\end{equation*}

局域展开的控制量不是单独由一个能量比猜测，而是上面已经定义的算符范数
\begin{equation*}
 \epsilon_{\rm Schur}=\sup_{p\in\mathcal D}\|K_M^{-1}\delta K(p)\|.
\end{equation*}
给定允许的二次核绝对误差 $\delta_K^{\rm tol}>0$，$N$ 阶局域 EFT 在目标能区 $\mathcal D$ 内满足精度要求的充分条件为
\begin{equation*}
 \|K_{LH}\|_{\mathcal D}\,\|K_{HL}\|_{\mathcal D}\,
 \|K_M^{-1}\|
 \frac{\epsilon_{\rm Schur}^{N+1}}{1-\epsilon_{\rm Schur}}
 \le \delta_K^{\rm tol},
 \qquad \epsilon_{\rm Schur}<1,
\end{equation*}
其中 $\|K_{LH}\|_{\mathcal D}\equiv\sup_{p\in\mathcal D}\|K_{LH}(p)\|$，$K_{HL}$ 同理。具体模型中外部能量与重质量之比可以用来估计 $\epsilon_{\rm Schur}$，但可接受性最终由这一算符余项界判定。Schur 补本身始终是精确代数恒等式；截断 Neumann 级数才定义局域 EFT 子模型。
\end{derivation}
